%%%%% MTecknology's Resume
%%%%
% Author: Michael Lustfield
% License: CC-BY-4
% - https://creativecommons.org/licenses/by/4.0/legalcode.txt
%%%%

\documentclass[a4paper,10pt]{article}

%%%%%%%%%%%%%%%%%%%%%%%
%% BEGIN_FILE: mteck.sty
%% NOTE: Everything between here and END_FILE can
%% be relocated to "mteck.sty" and then included with:
% \usepackage{mteck}

% Dependencies
% NOTE: Some packages (lastpage, hyperref) require second build!
\usepackage[empty]{fullpage}
\usepackage{titlesec}
\usepackage{enumitem}
\usepackage[hidelinks]{hyperref}
\usepackage{fancyhdr}
\usepackage{fontawesome5}
\usepackage{multicol}
\usepackage{bookmark}
\usepackage{lastpage}
\usepackage[T1]{fontenc}
\usepackage[utf8]{inputenc}
\usepackage[brazil]{babel}

% Sans-Serif
% \usepackage[sfdefault]{FiraSans}
% \usepackage[sfdefault]{roboto}
% \usepackage[sfdefault]{noto-sans}
% \usepackage[default]{sourcesanspro}

% Serif
\usepackage{CormorantGaramond}
\usepackage{charter}

% Colors
% Use with \color{Name}
% Can wrap entire heading with {\color{accent} [...] }
\usepackage{xcolor}
\definecolor{accentTitle}{HTML}{0e6e55}
\definecolor{accentText}{HTML}{0e6e55}
\definecolor{accentLine}{HTML}{a16f0b}

% Misc. Options
\pagestyle{fancy}
\fancyhf{}
\fancyfoot{}
\renewcommand{\headrulewidth}{0pt}
\renewcommand{\footrulewidth}{0pt}
\urlstyle{same}

% Adjust Margins
\addtolength{\oddsidemargin}{-0.7in}
\addtolength{\evensidemargin}{-0.5in}
\addtolength{\textwidth}{1.19in}
\addtolength{\topmargin}{-0.7in}
\addtolength{\textheight}{1.4in}
\setlength{\multicolsep}{-3.0pt}
\setlength{\columnsep}{-1pt}
\setlength{\tabcolsep}{0pt}
\setlength{\footskip}{3.7pt}
\raggedbottom
\raggedright

% ATS Readability
\input{glyphtounicode}
\pdfgentounicode=1

%-----------------%
% Custom Commands %
%-----------------%

% Centered title along with subtitle between HR
% Usage: \documentTitle{Name}{Details}
\newcommand{\documentTitle}[2]{
\begin{center}
{\Huge\color{accentTitle} #1}
\vspace{10pt}
{\color{accentLine} \hrule}
\vspace{2pt}
% {\footnotesize\color{accentTitle} #2}
\footnotesize{#2}
\vspace{2pt}
{\color{accentLine} \hrule}
\end{center}
}

% Create a footer with correct padding
% Usage: \documentFooter{\thepage of X}
\newcommand{\documentFooter}[1]{
\setlength{\footskip}{10.25pt}
\fancyfoot[C]{\footnotesize #1}
}

% Simple wrapper to set up page numbering
% Usage: \numberedPages
% WARNING: Must run pdflatex twice!
\newcommand{\numberedPages}{
\documentFooter{\thepage/\pageref{LastPage}}
}

% Section heading with horizontal rule
% Usage: \section{Title}
\titleformat{\section}{
\vspace{-5pt}
\color{accentText}
\raggedright\large\bfseries
}{}{0em}{}[\color{accentLine}\titlerule]

% Section heading with separator and content on same line
% Usage: \tinysection{Title}
\newcommand{\tinysection}[1]{
\phantomsection
\addcontentsline{toc}{section}{#1}
{\large{\bfseries\color{accentText}#1} {\color{accentLine} |}}
}

% Indented line with arguments left/right aligned in document
% Usage: \heading{Left}{Right}
\newcommand{\heading}[2]{
\hspace{10pt}#1\hfill#2\\
}

% Adds \textbf to \heading
\newcommand{\headingBf}[2]{
\heading{\textbf{#1}}{\textbf{#2}}
}

% Adds \textit to \heading
\newcommand{\headingIt}[2]{
\heading{\textit{#1}}{\textit{#2}}
}

% Template for itemized lists
% Usage: \begin{resume_list} [items] \end{resume_list}
\newenvironment{resume_list}{
\vspace{-7pt}
\begin{itemize}[itemsep=-2px, parsep=1pt, leftmargin=30pt]
}{
\end{itemize}
% \vspace{-2pt}
}

% Formats an item to use as an itemized title
% Usage: \itemTitle{Title}
\newcommand{\itemTitle}[1]{
\item[] \underline{#1}\vspace{4pt}
}

% Bullets used in itemized lists
\renewcommand\labelitemi{--}

%% END_FILE: mteck.sty
%%%%%%%%%%%%%%%%%%%%%%

%========================%
% Pedro Spegiorin Resume %
%========================%

% \numberedPages % NOTE: lastpage requires a second build
% \documentFooter{\thepage of 2} % Does similar without using lastpage

\begin{document}

%---------%
% Heading %
%---------%

\documentTitle{Pedro Spegiorin}{
% Web Version
% \raisebox{-0.05\height} \faPhone\ [redacted - web copy] ~
% \raisebox{-0.15\height} \faEnvelope\ [redacted - web copy] ~
%
\href{tel:55999118330}{
\raisebox{-0.05\height} \faPhone\ (55) 49 99911-8330} ~ | ~
\href{mailto:pedrospegiorindev36@gmail.com}{
\raisebox{-0.15\height} \faEnvelope\ pedrospegiorindev36@gmail.com} ~ | ~
\href{https://github.com/thehouseisonfire}{
\raisebox{-0.15\height} \faGithub\ github.com/thehouseisonfire}
}

%--------%
% Perfil %
%--------%

\tinysection{Quem Sou Eu}
Desenvolvedor Full Stack, mantenedor de projetos Open Source, estudante de Ciência da Computação e programador competitivo.

%--------%
% Skills %
%--------%

\section{Habilidades}

% \begin{multicols}{2}
\begin{itemize}[itemsep=-2px, parsep=1pt, leftmargin=105pt]
\item[\textbf{Idiomas}] Inglês (fluente)
\item[\textbf{Linguagens}] Rust, JavaScript/TypeScript, C++, Python, SQL, Bash
\item[\textbf{Frameworks}] React, Next.js, TanStack, Vue.js, Nuxt
\item[\textbf{Sistemas Operacionais}] Arch Linux, Ubuntu, Linux Mint
\item[\textbf{Ferramentas}] Git, GitHub CI/CD, Docker, Docker Compose, Supabase, Firebase, Vim
\end{itemize}
% \end{multicols}

%------------%
% Experience %
%------------%

\section{Experiência}

\headingBf{Desenvolvedor Júnior -- Difrisul Distribuidora Chapecó}{Fev 2025 -- Jul 2026}
\begin{resume_list}
\item Desenvolvi uma aplicação de e-commerce para a distribuidora utilizando Next.js, Supabase, Medusa.js, shadcn/ui, Tailwind CSS v4 e a SDK do Mercado Pago. Também iniciei o desenvolvimento de um aplicativo mobile integrado à plataforma.
\item Desenvolvi um script em Python para junção de rotas a partir de arquivos Excel, com geração de executável para Windows.
\item Desenvolvi um script em Python para filtragem e validação de números de telefone a partir de arquivos Excel, com geração de executável para Windows.
\end{resume_list}

\headingBf{Desenvolvedor Júnior -- Weplan Móveis Planejados}{Ago 2023 -- Fev 2025}
\begin{resume_list}
\item Desenvolvi atualizações para o site da empresa em produção, renovando seu estilo e conteúdo.
\item Atualizei toda a stack do site, incluindo a migração de Vue 2 para Nuxt 3.
\item Desenvolvi um novo site do zero utilizando Tailwind CSS e Nuxt 3.
\item Implementei diversas funcionalidades, correções e políticas para a ferramenta interna de cálculo e geração de contratos da empresa. Entre elas, persistência de dados entre telas, cálculo inverso a partir de um valor final, validação de dados, engenharia reversa de artefatos do software CAD PROMOB e extração de dados diretamente do XML resultante, suporte ao cálculo de múltiplos ambientes por contrato, otimização de aproximadamente 10x na criação de contratos, concatenação de PDFs e correção de diversos bugs e erros lógicos.
\item Atualizei toda a stack da ferramenta interna, incluindo a migração de Vue 2 para Nuxt 3, uma nova versão do Vuetify e atualizações das bibliotecas de desenvolvimento do Firebase.
\end{resume_list}

%------------%
% Activities %
%------------%

\section{Atividades}

\headingBf{Programação Competitiva}{2023 -- Presente}
\begin{resume_list}
\item Maratona de Programação (ICPC): medalha de bronze em Chapecó (2023--2024), medalha de bronze em Passo Fundo (2024--2025), medalha de prata em Chapecó (2025--2026) e medalha de ouro em Pinhalzinho (2026--2027).
\item Desbravalley Chapecó: 1º lugar (2023), 3º lugar (2024) e 2º lugar (2025).
\end{resume_list}

\headingBf{Monitoria}{Ago 2025 -- Dez 2025}
\begin{resume_list}
\item Algoritmos de Programação (2025.2)
\item Estruturas de Dados (2025.2)
\end{resume_list}

%-----------%
% Education %
%-----------%

\section{Formação Acadêmica}

\headingBf{Universidade Federal da Fronteira Sul}{2021 -- 2026}
\headingIt{Bacharelado em Ciência da Computação -- último semestre}{}

%-----------------%
% Recent Projects %
%-----------------%

\section{Projetos Recentes}

\headingBf{Rumqttc-next}{
\textcolor{blue}{
\href{https://github.com/thehouseisonfire/rumqtt/}{Link}
}
}
\begin{resume_list}
\item Biblioteca cliente MQTT completa, baseada em um fork criado após a empresa original cessar o desenvolvimento aberto do projeto.
\item Escrita em Rust, utilizando o runtime Tokio e Flume para comunicação MPSC.
\item Projeto amplamente utilizado, com centenas de milhares de downloads.
\end{resume_list}

\headingBf{Grok \& Z.ai Scraper}{
\textcolor{blue}{
\href{https://github.com/thehouseisonfire/grok-scraper}{Grok} /
\href{https://github.com/thehouseisonfire/zai-scraper}{Z.ai}
}
}
\begin{resume_list}
\item Scrapers simples para os sites grok.com e chat.z.ai.
\item Desenvolvidos originalmente em Python e posteriormente migrados para TypeScript.
\end{resume_list}

\headingBf{Image Tagger}{
\textcolor{blue}{
\href{https://github.com/thehouseisonfire/image-tagger}{Link}
}
}
\begin{resume_list}
\item Ferramenta para gerar listas de tags a partir de uma imagem, com suporte a quatro modelos de visão computacional.
\item Desenvolvida originalmente em Python e posteriormente migrada para TypeScript.
\end{resume_list}

\end{document}
